\section{Experimental Setup}
\label{sec:experimental-setup}

\subsection{Datasets}
\label{sec:datasets}

\textbf{MovieLens.} \texttt{ml-latest-small} \cite{movielens}: 100{,}836
ratings (0.5--5.0 stars, timestamped) from 610 users on 9{,}724 movies.
Item text is the title and genre tags (e.g.\ ``Toy Story (1995) Adventure
Animation Children Comedy Fantasy''), which is short, curated metadata.

\textbf{Amazon Reviews.} The \texttt{All\_Beauty} category of Amazon
Reviews 2023 \cite{amazon-reviews-2023}, restricted to its standard
5-core subset: 253 users, 356 items, and 2{,}535 ratings (1--5 stars).
Each item's text is its product title, its category, and up to 10 of its
reviews (title and body, each truncated to 500 characters).

\textbf{Yelp.} Businesses in Boise, ID, from the Yelp Open Dataset, with
at least 5 reviews: 2{,}937 businesses and 105{,}366 ratings (1--5
stars). 10{,}537 users have at least two positives and are evaluated
(median 3 positives). We chose a single mid-sized city to keep the
LLM embedding pass tractable on one machine. Item text is the business
name, its categories, and up to 10 reviews (500 characters each).

\textbf{Leak-free review text.} Review text is strongly correlated with
the star rating that defines relevance. If a held-out user's review were
part of an item's text, the label would leak into the features. We
therefore build item text only from reviews written by users who are
\emph{never evaluated}: reviewers outside the 5-core user set on Amazon,
and users with fewer than two positives on Yelp. No review whose rating
is a test label, under either split, can enter any item's text. After
this restriction the mean item text is 2{,}083 characters on Amazon and
2{,}148 on Yelp. Items with no qualifying review (18 of 356 on Amazon,
99 of 2{,}937 on Yelp) fall back to title and category.

\textbf{Within-dataset redundancy variants.} To test the redundancy
diagnostic (\S\ref{sec:redundancy-analysis}) without changing users, items, or
labels, we also run MovieLens with title-only and genres-only item
text.

\subsection{Task and candidate generation}
\label{sec:candidate-generation}
A rating $\geq 4.0$ is a positive. For each user with at least two
positives, we hold out 20\% of their positives (at least one) for test.
Each training positive is paired with 4 sampled negatives, and each test
list adds 50 sampled negatives drawn from items the user never rated.
Relevance is the rating for positives and 0 for negatives.

\textbf{Two negative-sampling regimes.} Sampled-negative metrics are
most lenient when negatives are easy \cite{krichene-sampled}.
Uniformly sampled unrated items are mostly obscure, so popularity alone
separates them from positives. We therefore report two regimes:
\emph{uniform} negatives, and \emph{popularity-matched} negatives drawn
with probability proportional to each item's rating count, which
removes most of the popularity shortcut. We treat popularity-matched
negatives as the primary setting.

\textbf{Splits.} The \emph{random} split shuffles each user's positives
before holding out 20\%. The \emph{temporal} split holds out each
user's most recent 20\%, so that the model must extrapolate to future
preferences. Popularity features are computed from every rating except
the held-out test positives. Under the temporal split they can
therefore include other users' later ratings, which is a mild,
pipeline-independent form of look-ahead that we note but do not remove.

\subsection{Metrics and statistics}
\label{sec:evaluation-metrics}
For each user we compute NDCG@10 with gain $2^{r}-1$ \cite{ndcg} and
Recall@10. Every configuration runs 5 seeds, and each seed redraws the
split and the negatives. Tables report the mean over seeds of the
per-user average, with a 95\% CI across seeds.

\textbf{Paired tests.} All pipelines see identical users, splits, and
candidate lists in each seed, so every comparison is paired. For each
pipeline pair we average each user's NDCG@10 over seeds and test the
per-user differences with (i) a paired $t$-test for any difference and
(ii) TOST equivalence \cite{tost} against a margin of $\pm0.01$
NDCG@10. The margin was fixed before the main experiments were run; it
equals half of the largest encoder gap observed in pilot runs with the
item-only scorer. Both families are Holm-corrected within each
scorer and configuration. We call a pair \emph{different} if only the
$t$-test rejects, \emph{equivalent} if only TOST rejects,
\emph{negligible} if both reject (a detectable difference inside the
margin), and \emph{inconclusive} if neither rejects.

\textbf{Latency.} Embedding latency is measured in a separate
benchmark, not inside the experiment runs. All four neural encoders
embed the same 256-item sample, with full-text coverage (chunking where
needed), at batch sizes 1 and 32, on the same machine with nothing else
running. Each model gets one warmup call, and three rounds rotate the
model order. We report the median per-item latency over rounds.
SBERT runs on the Metal (MPS) backend, mxbai through Ollama, and the
two GGUF models through llama.cpp servers with four parallel slots,
so a batch request is spread across slots. Every
pipeline runs locally, so the marginal API cost is \$0 and latency is
the cost signal.
